% ATTEMPT_STATUS: unresolved
% RESEARCH_GENERATED: 2026-10-01; GPT-6 Astra Ultra; individual mathematical attempt
% TOP_PROBLEM: {"id": 158, "title": "Sets with No Unique Sum Representations", "category_id": 2, "category": {"id": 2, "name": "combinatorics", "display_name": "Combinatorics", "description": "Counting problems, graph theory, discrete structures.", "slug": "combinatorics", "order_index": 2, "created_at": "2026-07-31T15:26:25.671Z"}, "status": "open"}
% FIRST_POSTED: 2026-10-01
% Imported from: unsolvedmath_additional_500.tex; UnsolvedMath ID 158
% !TeX program = lualatex
% Compile twice: lualatex 158.tex
% Self-contained: no external bibliography, images, or \input files.
% Numeric section labels and visible numbers are original UnsolvedMath IDs.
\documentclass[11pt,a4paper]{article}
\usepackage[margin=24mm,headheight=14pt]{geometry}
\usepackage{fontspec}
\setmainfont{Latin Modern Roman}
\setsansfont{Latin Modern Sans}
\setmonofont{Latin Modern Mono}
\usepackage{amsmath,amssymb,mathtools}
\usepackage{unicode-math}
\setmathfont{Latin Modern Math}
\usepackage{microtype}
\usepackage{enumitem,xurl,needspace}
\usepackage[dvipsnames]{xcolor}
\usepackage{hyperref}
\hypersetup{unicode=true,colorlinks=true,linkcolor=MidnightBlue,urlcolor=MidnightBlue,
 pdftitle={UnsolvedMath Problem 158},pdfauthor={Source-annotated compilation},bookmarksnumbered=true}
\usepackage{bookmark,tocloft,titlesec,fancyhdr}
\setlength{\cftsecnumwidth}{4.2em}
\cftsetpnumwidth{2.5em}
\cftsetrmarg{4em}
\titleformat{\section}{\Large\bfseries\raggedright}{\thesection}{0.7em}{}
\pagestyle{fancy}\fancyhf{}
\fancyhead[L]{\small\sffamily UnsolvedMath research attempt}
\fancyhead[R]{\small Original catalogue IDs}
\fancyfoot[C]{\thepage}
\setlength{\parindent}{0pt}
\setlength{\parskip}{5pt plus 1pt}
\setlength{\emergencystretch}{3em}
\setcounter{tocdepth}{1}\setcounter{secnumdepth}{1}
\setlist[itemize]{leftmargin=3em,itemsep=4pt,topsep=4pt,parsep=0pt}
\allowdisplaybreaks
\newcommand{\N}{\mathbb{N}}
\newcommand{\Z}{\mathbb{Z}}
\newcommand{\Q}{\mathbb{Q}}
\newcommand{\R}{\mathbb{R}}
\newcommand{\C}{\mathbb{C}}
\newcommand{\F}{\mathbb{F}}
\newcommand{\A}{\mathbb{A}}
\newcommand{\PP}{\mathbb{P}}
\newcommand{\E}{\mathbb{E}}
\newcommand{\eps}{\varepsilon}
\newcommand{\dd}{\,\mathrm d}
\newcommand{\sref}[2]{\hyperlink{source:#1:#2}{[S#2]}}
\newif\ifshowcatalogue
\showcataloguetrue
\newcommand{\cataloguescope}[1]{\ifshowcatalogue
 \paragraph{Statement recorded in the supplied source.}
 #1\fi}
\begin{document}
% UnsolvedMath ID 158; source catalogue code GREEN-070

\setcounter{section}{157}

\section{Small Sets without Uniquely Represented Sums}\label{158}

{\small\sffamily UnsolvedMath ID 158 · Combinatorics}

\subsection{Definitions and mathematical statement}

\paragraph{Shared notation.}Unless a section says otherwise, $\N=\{1,2,\ldots\}$,
$\Z$ is the ring of integers, $\Q,\R,\C$ are the rational, real, and complex
fields, $[N]=\{1,\ldots,\lfloor N\rfloor\}$, and $\log$ is the natural
logarithm. For real functions of a parameter tending to infinity,
$f\ll g$ and $f=O(g)$ mean $|f|\leq Cg$ eventually for some fixed $C$;
$f\gg g$ reverses this comparison; $f=o(g)$ means $f/g\to0$;
$f\sim g$ means $f/g\to1$; and $f\asymp g$ means both order bounds.
A subscript on $O$ or $\ll$ specifies allowed constant dependence.
The expression $n^{o(1)}$ in an upper bound means growth smaller than
$n^\epsilon$ for each fixed $\epsilon>0$. Local definitions govern any
exception, finite-field convention, or use of the same symbol for another
object. An asymptotic claim is not automatically an effective algorithm.



For $A\subseteq\Z/p\Z$, representations of $s$ are unordered pairs with repetition allowed: $(a,b)$ and $(b,a)$ count as one. Define $r_A^{\rm un}(s)=\#\{\{a,b\}_{\rm mult}:a,b\in A,\ a+b=s\}$. The condition is $r_A^{\rm un}(s)\ne1$ for every $s$; values outside $A+A$ have count zero. Minimize $|A|\geq2$ as a function of the prime $p$. \sref{158}{1} \sref{158}{2}

\paragraph{Statement recorded in the supplied source.}
What is the size of the smallest set $A \subset \mathbb{Z}/p\mathbb{Z}$ (with at least two elements) for which no element in the sumset $A + A$ has a unique representation? \sref{158}{1}

\paragraph{Scope and interpretation.}

The unordered-pair convention is specified in the source record’s cited Bedert paper and review. Counting the two orders as different would change the question. \sref{158}{1}

\paragraph{Residual task reported by the archive.}

The embedded review (2026-08-17) records: Close the gap between the super-logarithmic lower bound and quadratic-logarithmic upper bound, ideally determining the order or an asymptotic formula. \sref{158}{1}

\subsection{Short English statement}

How small can a set modulo a prime be if every sum it represents has at least two genuinely different representations?

\subsection{Research attempt}

\paragraph{Provenance and scope.}
This individual attempt was generated by GPT-6 Astra Ultra on 1 October 2026. The method was to derive explicit partial results and test the first proposed extension adversarially. The original statement and sources below are retained. No claim of mathematical novelty or complete solution is made.

\paragraph{Formal target and primary-source context.}
Representations are unordered pairs with repetition, so $(a,b)$ and $(b,a)$ count only once. We seek the smallest $|A|\geq2$ in $\Z/p\Z$ for which every represented sum has at least two such representations. Bedert's inspected paper studies substantially stronger asymptotic unique-sum bounds than the elementary argument below. We reconstruct a logarithmic rectification bound with explicit constants and give complete small-prime searches under the correct unordered convention.

\paragraph{Rectification by simultaneous approximation.}
Write $A=\{a_0,a_1,\ldots,a_{m-1}\}$ and suppose $4^{m-1}<p$. Consider the $4^{m-1}+1$ vectors
\[
 \left(\left\{\frac{t(a_i-a_0)}p\right\}\right)_{i=1}^{m-1},
 \qquad 0\leq t\leq4^{m-1},
\]
in the unit $(m-1)$-cube. Partition the cube into $4^{m-1}$ half-open subcubes of side one quarter. Two vectors fall in the same subcube. Subtracting their indices gives a nonzero integer $d$ with $|d|<p$ such that each $d(a_i-a_0)$ has a representative in $(-p/4,p/4)$ modulo $p$. Since $p$ is prime, multiplication by $d$ is invertible. Translation by $a_0$ and dilation by $d$ preserve representation multiplicities.

The transformed set therefore has integer representatives contained in an interval of length strictly less than $p/2$. Its integer pairwise sums lie in an interval of length less than $p$, so reduction modulo $p$ is injective on that sum range. Let $b$ be the smallest representative. The integer sum $2b$ has the unique unordered representation $\{b,b\}$: any other selected element is larger and raises the sum. By injectivity on the sum range, the same uniqueness holds modulo $p$.

Consequently any set with no unique sum must satisfy $4^{m-1}\geq p$, giving the explicit lower bound $m\geq1+\log_4p$ interpreted with integer rounding. This is only a basic logarithmic bound, not the stronger current lower bound from the cited research.

\paragraph{Why the interval-width condition is exact enough.}
The simultaneous approximation produces signed representatives on both sides of zero. It is their total span, less than $p/2$, that matters; a statement that each representative is merely less than $p/2$ in absolute value would give span almost $p$ and fail to prevent sum wraparound. Likewise, choosing two equal indices would give a zero dilation, which is why distinct pigeonhole points and the strict inequality $4^{m-1}<p$ are explicitly retained. The prime assumption supplies invertibility; an arbitrary composite modulus would need a coprimality argument absent here.

\paragraph{An elementary upper construction and finite minima.}
If $A\subset\F_p$ has size $m$, every $s$ has at least $2m-p$ ordered representations, because $A\cap(s-A)$ has that many points. Each unordered pair accounts for at most two ordered representations. Thus if $2m-p\geq3$, every sum has at least two unordered representations, and indeed every residue is represented. For odd $p$, any set of size $(p+3)/2$ supplies such an example. This linear-size upper bound is intentionally crude; it is included to check the counting convention and to provide a certified search endpoint.

The companion script searches all sets containing zero, which is sufficient because translation preserves the property. It tests increasing cardinalities and counts unordered pairs by combinations with repetition. The exact minima for $p=3,5,7,11,13$ are respectively $3,4,5,7,7$. For $p=13$, one witness is $\{0,1,2,3,6,8,10\}$; its representation counts are two at every residue except one residue with count four. The total is $28=7\cdot8/2$, verifying that no ordered-pair convention has slipped in.

\paragraph{Verification and remaining gap.}
The search is exhaustive for the stated finite primes, not a claim about larger ones. The logarithmic proof is independent of those computations and includes diagonal pairs, which provide its unique extreme sum. The basic upper construction works at $p=3$ as the whole group. The established output is a fully explicit rectification theorem and exact small cases. The unresolved asymptotic problem requires either a substantially stronger obstruction to sets that cannot be rectified into a half-length interval, or smaller verified constructions beyond the elementary dense-set argument. No new asymptotic record is claimed.

\subsection{Sources}

{\small

\begin{itemize}

\item[\hypertarget{source:158:1}{S1}] ULAM.AI, \emph{UnsolvedMath}, supplied \texttt{problems.json}, record 158 (GREEN-070), statement and background. The embedded literature-review date is 2026-08-17; that is the archive’s review date, not a fresh certification. Dataset landing page: \url{https://huggingface.co/datasets/ulamai/UnsolvedMath}.

\item[\hypertarget{source:158:2}{S2}] Benjamin Bedert, On unique sums in Abelian groups, Combinatorica 44 (2024), no. 2, 269-298; arXiv:2303.15134. \url{https://arxiv.org/abs/2303.15134}. Bibliographic information imported from the supplied record.

\item[\hypertarget{source:158:3}{S3}] Ben Green, 100 Open Problems, Problem 27, most recent update December 2025. \url{https://people.maths.ox.ac.uk/greenbj/papers/open-problems.pdf}. The author-hosted document was inspected on 27 September 2026; the archive’s local code is not a problem-number locator in this paper.

\item[E1] B. Bedert, \emph{On unique sums in Abelian groups}, current manuscript. \url{https://arxiv.org/abs/2303.15134}. Inspected 1 October 2026.

\end{itemize}

}

\label{end:158}
\end{document}
